%% kapa: exam=1 round=31 year=2020 period=1 form=A subject=real_estate subject_ko=부동산학원론 seq=18 number=98
%% source: Q-Net 감정평가사 기출문제 (cst00309) artlSeq=5207219 file=extracted/1008/5207219/A형/A형/2020년도 제31회 감정평가사 1차 시험_1교시_A형.pdf
%% conversion: kordoc 4.24.51 -> tools/gen1.py (auto); review=proofread-sonnet; flags=none
\begin{question}{98}
다음의 개발방식은?
\begin{framedtext}
○ 대지로서의 효용증진과 공공시설의 정비를 목적으로 하며, 택지개발사업에 주로 활용되는 방식이다. \par
○ 사업 후 개발토지 중 사업에 소요된 비용과 공공용지를 제외한 토지를 당초의 토지소유자에게 되돌려 주는 방식이다. \par
○ 개발사업 시 사업재원으로 확보해 놓은 토지를 체비지라고 한다. \par
\end{framedtext}
\begin{choices}
  \item 환지방식
  \item 신탁방식
  \item 수용방식
  \item 매수방식
  \item 합동방식
\end{choices}
\end{question}
\answer{1}
